\section{Problem Formulation}
Let the original training set be $D=D_r\cup D_f$, where $D_f$ is the requested deletion
set, $D_r$ is the retained set, and $D_r\cap D_f=\varnothing$. Let $A$ denote a randomized
training algorithm and $U$ an approximate unlearning operator. The original, exact, and
approximately unlearned models are
\begin{align}
M_0 &= A(D;s_0),\\
M_R^{(s)} &= A(D_r;s),\\
M_U &= U(M_0,D_f,D_r;\eta),
\end{align}
where $s$ indexes training randomness and $\eta$ contains method-specific randomness and
hyperparameters. The exact counterfactual is consequently the population
\begin{equation}
\mathcal{R}(D_f)=\{M_R^{(s)}:s\sim\mathcal{S}\},
\end{equation}
not a single model or parameter vector. P001 approximates this population with 15 fixed,
independent seeds.

For a probe set $X$ and discrepancy functional $d_j$, exact/exact variation is represented
by
\begin{equation}
\mathcal{N}_j=\{d_j(M_R^{(a)},M_R^{(b)};X):a<b\}.
\end{equation}
This set is an empirical null reference, but its $\binom{k}{2}$ entries are not independent
because each trained model appears in multiple pairs. Seed-level or manifest-level
resampling is therefore required for inferential procedures.

The original/exact distribution
\begin{equation}
\mathcal{O}_j=\{d_j(M_0,M_R^{(s)};X):s\in S\}
\end{equation}
tests whether deletion produces an observable signal relative to ordinary retraining
variation. Only after this reference-sensitivity condition is met do we interpret the
approximate/exact distribution
\begin{equation}
\mathcal{U}_j=\{d_j(M_U,M_R^{(s)};X):s\in S\}.
\end{equation}

Counterfactual fidelity is kept multidimensional. A claim can concern predictions,
probability distributions, calibration, privacy attacks, intermediate representations,
local neighborhoods, robustness under controlled corruptions, retained utility, or
computational cost. Agreement on one coordinate does not imply agreement on another.
